\section{Pilot}
\label{sec:Pilot}
\subsection{Participants}
A total of 61 undergraduate students from Ontario Tech University participated in the pilot study as part of the SONA participant pool.
This sample consisted of 61 participants (20 males and 41 females, M = 20.84, SD = 3.25, range = 17 to 38).

\subsection{Stimuli}
\label{sec:Stimuli}
One hundred and forty images of facial expressions from the racially diverse affective expression (RADIATE) and UIBVFED datasets were used \citep{conley_racially_2018, oliver_uibvfed_2020}.
The RADIATE contains perceptually validated images of racially and ethnically diverse participants, aged 18-30 years old, each expressing 16 different emotions. 
The UIBVFED dataset contains a set of 20 virtual characters that are also ethnically diverse, aged 20-80 years old expressing 32 emotions. 
The UIBVFED facial expressions were created using blendshapes, a tool that represents and manipulates clusters of facial landmarks similar to that of facial action units (AUs).
Then adult models (5 males and 5 females) from each dataset were identified and matched between-sets on face shape, sex, skin tone, and hair colour. 
Images of each model expressing seven emotions (Anger, Disgust, Fear, Joy, Sadness, Surprise, Neutral) were selected. 
Expressions were selected for each model, that closely align with Ekman's 6 basic emotions + neutral \citep{ekman1971constants}.
The intensity of the emotional expressions were matched as well across both datasets. 

\subsection{Design}
The first stage of the task required participants to select which emotion was portrayed by an individual virtual avatar.
For each trial, they were shown a fixation cross for one second, then a photo of an avatar expressing an emotion was displayed in the center of the screen.
Participants were then asked to identify the emotion being expressed by selecting one of the seven emotions. 
Only the virtual avatars from the UIBVFED stimulus set were used in this emotion recognition task, as the RADIATE dataset has already been validated for emotion recognition. 

The second stage of the task required participants to rate the similarity of expressions displayed by virtual avatars compared to real human models.
Participants were given instructions to rate the similarity of the expressions displayed on the faces, not of the faces themselves.
Participants were shown a fixation cross for one second.
Following this, a pair of photos was selected randomly and displayed side-by-side on screen, with a slider bar underneath to make the similarity rating.
There were three conditions in the similarity matching task: (1) a congruent-congruent (CC) condition, in which photos were matched on model similarity and expression similarity, (2) a congruent-incongruent (CI) condition, in which photos were matched on model similarity but not expression similarity, and (3) an incongruent-congruent (IC) condition, in which photos were not matched on model similarity but were matched on expression similarity.

\begin{figure}[H]
    \centering
    \begin{subfigure}[b]{0.49\textwidth}
        \centering
        \includegraphics[width=\linewidth]{C:/Users/super/OneDrive - Ontario Tech University/fNIRS_Emotions/plots/zedd_figures/emotion_recognition_task.png}
        \caption{Example trial from the emotion recognition task using virtual avatars from the UIBVFED stimulus set.}
        \label{fig:emotion_recognition_task}
    \end{subfigure}\hfill
    \begin{subfigure}[b]{0.49\textwidth}
        \centering
        \includegraphics[width=\linewidth]{C:/Users/super/OneDrive - Ontario Tech University/fNIRS_Emotions/plots/zedd_figures/similarity_task.png}
        \caption{Example trial from the similarity matching task (congruent-congruent) showing a virtual avatar and a real human model.}
        \label{fig:similarity_task}
    \end{subfigure}
    \caption{Left: emotion recognition trial. Right: similarity matching trial.}
    \label{fig:task_examples_side_by_side}
\end{figure}

\subsection{Results}
\subsubsection{Accuracy}
To calculate the effect of target emotion on accuracy of emotion recognition, each participant's average accuracy was calculated for each emotion category, as well as for overall accuracy. 
Mean and standard deviation of recognition accuracy for each emotion and dataset is reported in Table \ref{tab:accuracy comparison}.

\input{../../../processed_data/zedd_data/accuracy_comparison_table.tex}

\subsubsection{Similarity Ratings}
To assess whether the physical features of the faces pictured had an effect on the similarity ratings, pairwise $t$-tests were run comparing the CC and IC conditions using a Holm correction. 
The only emotion that did not show a significant difference in similarity rating regardless of the physical features of the face was Sadness ($t(12{,}999) = -0.051$, $p = 0.95$).

\begin{figure}[H]
    \centering
    \includegraphics[width=1\textwidth]{C:/Users/super/OneDrive - Ontario Tech University/fNIRS_Emotions/plots/zedd_figures/similarity_means_grouped_plot.png}
    \caption{Mean similarity ratings for each emotion across congruent-congruent (CC) and incongruent-congruent (IC) conditions. 
    The n.s. indicates that there was no significant difference between the two conditions for the Sadness emotion.}
    \label{fig:similarity_ratings}
\end{figure}